- When adding extra gateway, devices will automatically update SF (spreading factor)
- ADR works based on strength: devices closer to gateways use high data rate
- ADR only works for fixed devices, not mobile devices
- The Network Reference Model (NRM) is used to identify the interfaces between elements used in a LoRaWAN network
- RP002 is the companion document to TS001
- 3 factors to consider when deciding on a network:
    - price 
    - gateway access control 
    - privacy
- public networks:
    - verify network coverage before subscription
    - only makes sense if not too many devices
    - cant optimize by adding own gateways
- private networks:
    - for large number of devices 
    - huge overhead: gateways have to be installed and maintained

- The network server is where the LoRaWAN networking protocol terminates, and where data handling takes place
- One can use Semtech's network server for free when prototyping: up to 3 gateways and 10 end devices 
- The join server is used when an end device is joined to a network via Over-the-Air Activation (OTAA) or when a device requests 
  new security session keys via the Rejoin MAC command

- Each end dvice consists of:
    - MCU that runs LoRaWAN stack + power supply
    - peripherals, e.g sensors
    - LoRa Radio
    - Antenna
- Device software layers:
    - drivers
    - middleware
    - application 

- 1/4 wavelength for 868MHz = 8.2cm
- distinction between indoor and outdoor gateways

- LoRaMAC-Node from Semtech: free open-source project
- Also ST's LoRaWAN stack for all STM32Lxx MCUs
- Mbed stack for C++

- LoRaWAN always adds 13 bytes to application payload

- One could, for example, transmit a min|avg|max every five minutes, or you could transmit only when a sensor value changes more than a certain threshold amount. 
  Similarly, you could have transmissions be triggered by motion or other events.
- The data rate should be as fast as possible, to minimize your airtime. 
  Using SF7 and BW125 (SF7BW125) is usually a good place to start.

- FUOTA: firmware update over the air

- Over-the-air-activation (OTAA) vs activation-by-personalization (ABP)
- ABP: device has all the info before connecting to a network + DevEUI + DevAddr
- OTAA: device must issue join request followed by joinaccept from network
- While provisioning devices with ABP, the best practice is to also provision 
  them with all the elements necessary for OTAA (i.e., DevEUI, JoinEUI, and Root Keys)
- OTAA is preferred approach
- All devices must keep a set of persistent values, even after power loss or reset

- Transmission is very power consuming (10 times more than receiving)

- To secure the end devices deployed in the field, crypto chips can be used. Crypto chips, commonly known as secure elements 
  can securely store keys and efficiently encrypt and decrypt data.
- In crypto chips: registers in secure element are write only
- Crypto chips can also be used to verify the authenticity of the software that is uploaded to the device. 
  This mitigates the threat of an attacker overwriting the device firmware
- crypto chip stores network key and derived session keys 
- fast encryption without burdening CPU 
- When using LoRaWAN, ensure that the join server runs in a secure domain that you host or that is hosted by a trusted third party, because this server stores the root keys for the cloud 
  side of the LoRaWAN network. To prevent vendor lock-ins, check to see whether the join server is 
  network-provider agnostic and whether it allows end devices to join a different LoRaWAN network operator over time.
- Out-of-band key provisioning can be a good way to make sure that the people who insert the keys don’t have access to the keys of the devices themselves. 
  It is possible, however, to create an application that uses BLE or NFC to inject the keys directly into the device.
- To mitigate against replay attacks, enable 32-bit frame counters for uplink and downlink messages. Also, use a fixed payload length to mask event-driven 
  activities and help prevent hackers from deriving event-specific information from the payload

- Signal strength testing vs radio planning
- Lookup how to calculate LoRa link budget
- 6dB loss for doubling distance
- FSPL: -150dB/800km
- lookup multipath fading, Fresnel Zone
- 5 to 10 dB of safety margin for link budget
- According to two ray model: -12dB @ 2x distance
- max antenna gain: +2.15dBi (=dipole gain)

- if the internal walls are not brick or concrete, it should be possible to cover most of the 
  building’s interior with a single gateway on the same floor as the end devices.
- You can use radio planning tools to improve the coverage of your network
- There are some tools on the market that have LoRa modules available, such as Forsk Atoll, Infovista Planet, or iBwave
- outdoor gateways are more expensive than indoor due to waterproofness, cavity filters etc
- for indoor: put gateways near glass
- MAPL: max allowable path loss
- Friis eq for free space path loss (FSPL): 20log... 
- uplink and downlink budget are not equal, due to different elevations
- Propagation models:
  1. Log path loss (allrounder)
  2. okumura: most used for 150-1500MHz
  3. Cost 231: goes into more details on city architecture
  4. ITU-R P.1238 (indoor)
- RF planning methodology:
  1. analyze link budget 
  2. RF coverage simulation
  3. field tests
- Tools (open source):
  1. Radio mobile, CloudRF (Web-based commercial product with free plan)
  2. SDR# combined with RTL-SDR for spectrum analysis
  3. GNU Radio 
  4. QGIS - geographic information system
- Radio planning test:
  1. send packets containing sequence number and GPS coordinates to keep track of packet loss \cite{petajajarviCoverageLPWANsRange2015}
  2. heat map can be generated using Google Maps JavaScript API % cite same as above
Radio interface system planning book
- microcell: anntenna height below rooftop level
- macro: (Radio Interface system planning book)
- important when radio planning: coverage, capacity, quality
- radio system planning phases:
  - dimensioning
  - detailed radio system planning 
  - optimisation
- diversities to combat fading:
  - space
  - time 
  - freq
  - polarization 

- Development stacks:
  - The things stack (needs internet)
  - ChirpStack (more offgrid promising), on RaspPi for example
  - zephyr with integrated LoRaWAN stack 
  - renode as hardware simulator
  
- Hardware to buy:
  - nordic power profiler kit II (PPK2)
  - RTL-SDR (software defined radio)
  - For gateway: Dragino LPS8 / MikroTik LR8 / DIY Rasp pi / Pycom LoPy4 / TTGo Lora 32
  - MCU: STM32WL55 (industrial favorite) / Heltech Wifi Lora 32 (ESP-based with OLED) / RakWireless WisBlock (for low-power)
  - Klassifikation MCUs: Prototypen vs produktionsfertig
  - STM32WL55: output transmit power up to 22dBm @V_DDPA=3.3V
  - STLink


- Software tools:
  - draw.io / Lucidchart
  - Node-Red for data visualization
  - MQTT for debugging
  - octopart for price comparison
  - semtech lora calculator
  - RF line of sight calculator

PCB design:
  - for printed PCB: look at reference designs
  - most notable: DN024 from TI: monopole 868MHz

% END NOTES

\subsection{Strategien zur Datenmengereduktion}
\label{subsec:datenreduktion}

Neben der Wahl des Spreading Factors lässt sich der Energieverbrauch durch die
Reduktion der Sendefrequenz und der Datenmenge weiter minimieren. Anstatt Rohdaten
kontinuierlich zu übertragen, können aggregierte Werte -- beispielsweise Minimum,
Durchschnitt und Maximum über ein Zeitfenster von fünf Minuten -- gesendet werden.
Alternativ können Übertragungen ereignisgesteuert erfolgen, etwa wenn ein Messwert
einen definierten Schwellwert überschreitet oder eine Bewegung detektiert wird.
\textit{Firmware Updates Over the Air} (FUOTA) ermöglicht darüber hinaus die
Aktualisierung der Gerätesoftware ohne physischen Zugang~\cite{semtech_lora_basics}.

% -----------------------------------------------------------------------------
\section{RF"=Planung und Signalausbreitung}
\label{sec:rf-planung}
% -----------------------------------------------------------------------------

RF"=Planung (\textit{Radio Frequency Planning}) bezeichnet die systematische
Vorgehensweise zur Vorhersage der Funkabdeckung eines drahtlosen Systems.
Sie bildet die methodische Grundlage für den in Kapitel~\ref{chap:evaluation}
durchgeführten Vergleich zwischen theoretisch vorhergesagter und praktisch
gemessener Reichweite.

% -----------------------------------------------------------------------------
\section{Energieeffizienz in eingebetteten LoRaWAN"=Systemen}
\label{sec:energieeffizienz}
% -----------------------------------------------------------------------------

Energieeffizienz ist eine der zentralen Anforderungen an das in dieser Arbeit
entwickelte System. Der Großteil des Energieverbrauchs entfällt auf den Sendevorgang:
Typischerweise verbraucht das Senden etwa zehnmal mehr Energie als der
Empfangsbetrieb~\cite{semtech_lora_basics}. Für den STM32WLE5 beträgt die Stromaufnahme
im Sendebetrieb bei maximaler Ausgangsleistung von \SI{22}{dBm} (bei $V_\text{DDPA} =
\SI{3.3}{\V}$) bis zu \SI{87}{\mA}, im Empfangsbetrieb typischerweise \SI{5.5}{\mA},
während der Deep"=Stop"=Modus lediglich \SI{1.08}{\uA} benötigt~\cite{stm32wl_ds}.

\subsection{Sleep"=Modi und Power Management}
\label{subsec:sleep}

Zephyr RTOS bietet ein umfassendes Power"=Management"=Framework, das es ermöglicht,
den Mikrocontroller zwischen Messzyklen in einen energiesparenden Schlafzustand zu
versetzen. Das \textit{Tickless Idle}"=Konzept deaktiviert den System"=Tick"=Timer
in Ruhephasen vollständig; einzelne Peripheriekomponenten können über
\texttt{PM\_DEVICE\_STATE\_SUSPEND} individuell abgeschaltet werden. Der STM32WLE5
unterstützt mehrere Schlafmodi mit unterschiedlichen Aufwachzeiten und
Stromverbräuchen, vom einfachen Sleep über Stop"=Modi bis hin zum Standby.

\subsection{Freiraum"=Pfadverlust und Ausbreitungsmodelle}
\label{subsec:pfadverlust}

Der Freiraum"=Pfadverlust (\textit{Free Space Path Loss}, FSPL) beschreibt den
Signalverlust in einer hindernisfreien Umgebung und berechnet sich nach der
Friis"=Gleichung:

\begin{equation}
  \text{FSPL} \, [\si{\dB}] = 20 \log_{10}(d) + 20 \log_{10}(f)
                              + 20 \log_{10}\!\left(\frac{4\pi}{c}\right)
  \label{eq:fspl}
\end{equation}

mit der Distanz $d$ in Metern und der Frequenz $f$ in Hertz. Bei Verdopplung der
Entfernung steigt der Pfadverlust im Freiraummodell um \SI{6}{\dB}; das
\textit{Two"=Ray Ground Reflection Model} ergibt bei Bodennähe sogar einen Verlust
von \SI{12}{\dB} bei doppelter Distanz~\cite{dobkin2012rf}. Für sehr große Distanzen
ergibt sich für das \SI{868}{\MHz}"=Band bei beispielsweise \SI{800}{\km} ein
theoretischer FSPL von \SI{-150}{\dB}.

Für realitätsnähere Vorhersagen werden empirische Ausbreitungsmodelle verwendet.
Tabelle~\ref{tab:ausbreitungsmodelle} gibt einen Überblick der gängigsten Modelle.

\begin{table}[htbp]
  \centering
  \caption{Übersicht der Ausbreitungsmodelle und deren Einsatzbereich}
  \label{tab:ausbreitungsmodelle}
  \begin{tabular}{@{}lll@{}}
    \toprule
    \textbf{Modell} & \textbf{Einsatzbereich} & \textbf{Bemerkung} \\
    \midrule
    Log"=Distance Path Loss & Allgemein & Einfachstes empirisches Modell \\
    Okumura"=Hata            & 150--1500\,MHz, Außenbereich & Meistgenutztes Modell \\
    COST 231                 & 1500--2000\,MHz              & Detaillierte Stadtarchitektur \\
    ITU"=R P.1238            & Innenbereich                 & Für Gebäudeplanung \\
    \bottomrule
  \end{tabular}
\end{table}

Das \textbf{Okumura"=Hata}"=Modell ist für den Frequenzbereich von 150 bis
\SI{1500}{\MHz} das am häufigsten eingesetzte Modell und eignet sich damit direkt für
das \SI{868}{\MHz}"=Band~\cite{haslett2008essentials}. Für Innenräume, wie sie
bei der Messung in Industriegebäuden auftreten, ist das \textbf{ITU"=R P.1238}"=Modell
geeigneter. Bei der Interpretation von Messergebnissen sind zudem
Mehrwegeausbreitung (\textit{Multipath Fading}) und die Fresnel"=Zone zu
berücksichtigen: Eine freie Fresnel"=Zone erfordert, dass das Gelände innerhalb einer
ellipsoidförmigen Zone um die direkte Sichtlinie frei von Hindernissen ist.

\subsection{RF"=Planungsmethodik}
\label{subsec:rf-methodik}

Die empfohlene Methodik zur RF"=Planung folgt einem dreistufigen Prozess~\cite{haslett2008essentials}:

\begin{enumerate}
  \item \textbf{Link"=Budget"=Analyse:} Berechnung des MAPL anhand der Systemparameter
        (Sendeleistung, Antennengewinn, Empfängerempfindlichkeit, Sicherheitsmarge).
  \item \textbf{RF"=Coverage"=Simulation:} Verwendung von Ausbreitungsmodellen und
        digitalem Geländemodell zur kartografischen Darstellung der vorhergesagten Abdeckung.
  \item \textbf{Feldmessungen:} Verifikation der Simulation durch reale Messungen in den
        Zielumgebungen.
\end{enumerate}

Für die Coverage"=Simulation steht mit \textbf{SPLAT!} ein quelloffenes Linux"=Werkzeug
zur Verfügung, das das Longley"=Rice"=Modell mit \acrfull{srtm}"=Geländedaten kombiniert.
Die gemessenen RSSI"=Werte können anschließend mit Python und der Bibliothek Folium
als georeferenzierte Heatmap über OpenStreetMap visualisiert und mit der
SPLAT!"=Simulation überlagert werden.

\subsection{Mikrozellen und Makrozellen}
\label{subsec:zellen}

In der Netzplanung wird zwischen \textit{Mikrozellen} -- bei denen die Antenne
unterhalb der Dachlinie installiert wird -- und \textit{Makrozellen} mit Antennen
oberhalb der Dachlinie unterschieden. Diese Unterscheidung beeinflusst das
gewählte Ausbreitungsmodell und die erreichbaren Abdeckungsradien~\cite{walter2006radio}.
Für Indoor"=Installationen gilt als Faustregel: Wenn die Innenwände nicht aus Beton oder
Mauerwerk bestehen, kann ein einzelnes Gateway auf der gleichen Etage wie die
Endgeräte ausreichende Abdeckung liefern.

% TODO: see where this section fits
\subsection{Öffentliche vs. private Netzwerke}
\label{subsec:netzwerktypen}

Bei der Wahl der Netzwerkinfrastruktur sind drei Faktoren maßgeblich: Kosten,
Zugangskontrolle und Datenschutz~\cite{LoRaModulationBasics}.

Öffentliche Netze (z.B. The Things Network) sind für kleine Gerätezahlen und
Prototyping geeignet, erfordern jedoch eine bestehende Abdeckung und bieten keine
Möglichkeit zur Optimierung durch eigene Gateways. Private Netzwerke eignen sich
für große Gerätezahlen und sicherheitskritische Anwendungen, erfordern jedoch die
Installation und den Betrieb eigener Gateways und Server"=Infrastruktur. Für das
in dieser Arbeit betrachtete Off"=Grid"=Szenario -- den Einsatz in abgelegenen
Gefahrgebieten ohne bestehende Infrastruktur -- ist ein privates, autarkes Netzwerk
die einzig praktikable Option. Als Network"=Server"=Software kommt dabei \textbf{ChirpStack}
zum Einsatz, da dieser vollständig lokal auf einem Raspberry Pi betrieben werden kann
und keine Internetverbindung erfordert. Semtech bietet für Prototyping"=Zwecke zudem
einen kostenlosen Network Server für bis zu drei Gateways und zehn Endgeräte an.





